% ============================================================================
% CAPÍTULO 1 — O universo full stack\index{full stack} e o mercado
% ============================================================================
\chapter{O Universo Full Stack e o Mercado}
\label{cap:universo}

\begin{caixaconceito}
\textbf{Full stack} é a capacidade de construir um produto web de ponta a
ponta: a interface que o usuário vê (frontend\index{frontend}), a lógica e os dados que rodam
no servidor (backend\index{backend}), o banco de dados, e ainda o ambiente de implantação
(devops/cloud) — com qualidade, segurança e performance.
\end{caixaconceito}

Ao final deste capítulo, você será capaz de:
\begin{objetivos}
  \item Explicar o que é um desenvolvedor full stack e por que essa carreira é valorizada;
  \item Descrever, com dados de mercado, as tecnologias mais usadas em 2025/2026;
  \item Justificar a stack escolhida para este livro (TypeScript de ponta a ponta);
  \item Explicar o caminho de uma requisição HTTP\index{HTTP}, do navegador até o servidor e de volta;
  \item Configurar as ferramentas essenciais: editor, terminal, DevTools e Git;
  \item Publicar seu primeiro site estático no GitHub Pages (projeto do capítulo).
\end{objetivos}

\section{O que é desenvolvimento full stack}

Imagine um restaurante. O \textbf{salão} é o que o cliente vê: o cardápio, a
decoração, o atendimento. A \textbf{cozinha} é onde a comida é preparada: os
processos, os ingredientes, os estoques. O \textbf{estoque e a contabilidade}
são os dados que mantêm o negócio de pé. Um restaurante só funciona bem quando
os três funcionam juntos — e quem entende do sistema inteiro consegue tomar
decisões melhores.

No desenvolvimento web, a analogia é direta:
\begin{itemize}[leftmargin=*, itemsep=0.3em]
  \item \textbf{Frontend} (o salão): HTML, CSS e JavaScript\index{JavaScript} executados no
  navegador. É tudo o que o usuário vê e com o que interage.
  \item \textbf{Backend} (a cozinha): servidores, APIs, regras de negócio e
  autenticação. É o código que roda fora do navegador, em um servidor.
  \item \textbf{Banco de dados} (o estoque): o lugar onde os dados vivem de
  forma organizada, confiável e consultável.
  \item \textbf{Devops e cloud} (a infraestrutura): como o sistema é
  empacotado, testado, implantado e monitorado em produção.
\end{itemize}

O desenvolvedor full stack transita por todas essas camadas. Ele não precisa
ser o melhor do mundo em cada uma — precisa \emph{entender o sistema inteiro} e
ser capaz de produzir valor em qualquer ponto dele.

\begin{caixadica}
O termo \emph{full stack} não significa ``fazer tudo sozinho''. Significa ter
visão de sistema: saber conversar com especialistas de frontend, backend,
dados e infraestrutura na mesma língua. É exatamente essa habilidade que as
empresas mais procuram em 2026.
\end{caixadica}

\section{O que os dados dizem: as tecnologias mais usadas}

Este livro não escolhe tecnologias por modismo. A escolha foi baseada em três
levantamentos independentes e públicos, consultados em agosto de 2026:

\subsection{Pesquisa Stack Overflow 2025}
A pesquisa anual da Stack Overflow recebeu mais de 49 mil respostas de 177
países, cobrindo 314 tecnologias \cite{so2025,so2025geral}. Os pontos centrais:
\begin{itemize}[leftmargin=*, itemsep=0.3em]
  \item \textbf{JavaScript} é a linguagem mais usada do planeta — no topo em
  praticamente todas as edições desde 2013;
  \item \textbf{React} é o framework web mais usado (aproximadamente 45\% dos
  desenvolvedores);
  \item \textbf{PostgreSQL} é o banco de dados nº 1 (aproximadamente 58\% de
  adoção), à frente de MySQL, SQLite e MongoDB;
  \item \textbf{Node.js} lidera entre os runtimes e ambientes de execução de
  JavaScript (aproximadamente 49\%);
  \item \textbf{SQL} permanece entre as três tecnologias mais populares em
  geral — banco de dados é habilidade obrigatória;
  \item \textbf{TypeScript} segue em ascensão constante ano após ano.
\end{itemize}

\subsection{Octoverse 2025 (GitHub)}
O relatório anual do GitHub documentou um marco histórico: em agosto de 2025,
\textbf{o TypeScript ultrapassou Python e JavaScript} e tornou-se a linguagem
mais usada do GitHub, a maior plataforma de código do mundo \cite{octoverse2025}.
A plataforma ganha um novo desenvolvedor a cada segundo (mais de 180 milhões de
desenvolvedores). O crescimento está fortemente ligado ao avanço da
Inteligência Artificial: grande parte do código gerado por assistentes de IA é
escrito em TypeScript/JavaScript — o que torna essa linguagem a mais importante
para quem quer ser produtivo com IA.

\subsection{State of JavaScript 2024}
O levantamento da comunidade \cite{stateofjs2024} confirma a mesma direção:
React\index{React} mantém a liderança entre os frameworks; Vue segue em segundo; e o
TypeScript se consolidou como padrão de facto para projetos sérios de
JavaScript.

\subsection{Síntese dos dados}
\begin{table}[h]
\centering
\small
\caption{As tecnologias mais usadas e a decisão editorial deste livro.
Fontes: \cite{so2025}, \cite{octoverse2025}, \cite{stateofjs2024}.}
\label{tab:mercado}
\begin{tabularx}{\textwidth}{@{}>{\hsize=0.8\hsize\raggedright\arraybackslash}X >{\hsize=1.2\hsize\raggedright\arraybackslash}X >{\hsize=1.0\hsize\raggedright\arraybackslash}X >{\hsize=1.0\hsize\raggedright\arraybackslash}X@{}}
\toprule
\textbf{Camada} & \textbf{Tecnologias líderes} & \textbf{Adoção} & \textbf{Escolha deste livro} \\
\midrule
Linguagem & JavaScript / TypeScript & JS nº 1 no SO Survey; TS nº 1 no GitHub & TypeScript \\
Frontend & React & \textasciitilde 45\% & React + Next.js \\
Framework full stack & Next.js & Meta-framework React mais usado & Next.js 16 \\
Estilo & Tailwind CSS & Padrão de mercado & Tailwind v4 \\
Backend & Node.js\index{Node.js} & \textasciitilde 49\% & Node.js 24 LTS \\
Banco de dados & PostgreSQL\index{PostgreSQL} & \textasciitilde 58\% & PostgreSQL 18 \\
ORM & Prisma / Drizzle & Líderes no ecossistema TS & Prisma e Drizzle \\
Contêineres & Docker & Padrão da indústria & Docker + Compose \\
CI/CD & GitHub Actions & Integração nativa com GitHub & GitHub Actions \\
Cloud & AWS & Provedor líder & AWS \\
Testes & Vitest + Playwright & Padrão moderno TS & Vitest + Playwright \\
\bottomrule
\end{tabularx}
\end{table}

\begin{caixaconceito}
\textbf{Por que TypeScript de ponta a ponta?} Uma única linguagem no frontend,
no backend, nos scripts de infraestrutura e nos testes reduz drasticamente a
carga cognitiva de quem está aprendendo. Você domina uma sintaxe, um sistema de
tipos e um ecossistema de pacotes — e constrói o produto inteiro com ela. É o
caminho com menor atrito entre ``entendi o conceito'' e ``entreguei o sistema''.
\end{caixaconceito}

\section{Como a web funciona: a jornada de uma requisição}

Antes de escrever código, é preciso entender o palco onde ele vai se apresentar.
Quando você digita \texttt{https://exemplo.com.br} no navegador, acontece o
caminho da Figura~\ref{fig:jornada-requisicao} — os passos a seguir detalham
cada etapa:

\begin{figure}[htbp]
  \centering
  \diagramatikz{%
  \begin{tikzpicture}[
      box/.style={draw=primaria, fill=azulclaro, rounded corners=2pt,
        align=center, inner sep=4pt, font=\footnotesize},
      seta/.style={-{Stealth[length=2.2mm]}, thick, color=secundaria},
      font=\footnotesize
    ]
    \node[box] (u) at (0,0) {Usuário};
    \node[box] (n) at (3.0,0) {Navegador};
    \node[box] (s) at (6.1,0) {Servidor};
    \node[box] (b) at (9.5,0) {Banco de\\dados};
    \draw[seta] (u) -- node[above] {1. DNS} (n);
    \draw[seta] (n) -- node[above] {2. HTTPS} (s);
    \draw[seta] (s) -- node[above] {3. SQL} (b);
    \draw[seta] (b) to[out=-70,in=-110] node[below] {4. resposta} (s);
    \draw[seta] (s) to[out=-70,in=-110] node[below] {5. HTML/CSS/JS} (n);
    \draw[seta] (n) to[out=-70,in=-110] node[below] {6. renderiza} (u);
  \end{tikzpicture}}
  \caption{A jornada de uma requisição: ida (1--3) e resposta (4--6).}
  \label{fig:jornada-requisicao}
\end{figure}

\begin{passos}
  \item O navegador pergunta ao \textbf{DNS} (Domain Name System) qual é o
  endereço IP do domínio \texttt{exemplo.com.br} — como consultar uma agenda
  telefônica para achar o número de um nome;
  \item Com o IP em mãos, o navegador abre uma conexão segura (TLS) com o
  servidor e envia uma \textbf{requisição HTTP} — tipicamente um \texttt{GET}
  pedindo o documento da página;
  \item O servidor processa a requisição\index{requisição} e responde com um \textbf{código de
  status} (\texttt{200 OK}, \texttt{404 Não encontrado}, etc.) e o conteúdo —
  geralmente HTML, CSS e JavaScript;
  \item O navegador \textbf{renderiza} o HTML, aplica o CSS e executa o
  JavaScript, montando a página que você vê;
  \item Se a página precisa de mais dados (por exemplo, a lista de produtos),
  o JavaScript faz novas requisições (via \texttt{fetch}) a uma \textbf{API},
  que lê o banco de dados e devolve JSON.
\end{passos}

\begin{caixaconceito}
\textbf{HTTP} (Hypertext Transfer Protocol) é o protocolo de comunicação da
web. Cada interação é um par requisição--resposta: o cliente pede, o servidor
responde com um status e um corpo. Você estudará isso a fundo no
capítulo~\ref{cap:http}.
\end{caixaconceito}

\begin{caixadica}
Para ver essa jornada acontecendo ao vivo: abra o \textbf{DevTools} do
navegador (F12), aba \emph{Rede/Network}, recarregue qualquer página e observe
as requisições. Cada linha é um pedido HTTP com seu status, tamanho e tempo de
resposta. Esse é o raio-X da web — e será seu melhor amigo em toda a carreira.
\end{caixadica}

\section{As ferramentas essenciais}

Um desenvolvedor full stack vive em cinco ferramentas:

\begin{enumerate}[leftmargin=*, itemsep=0.4em]
  \item \textbf{Editor de código} — o VS Code é o padrão de mercado
  (gratuito, extensível e com suporte nativo a TypeScript);
  \item \textbf{Terminal} — o ambiente de linha de comando onde você roda
  Git, Node, npm e Docker. No Windows, o Git Bash (instalado junto com o Git);
  \item \textbf{DevTools do navegador} — inspeção de HTML/CSS, console de
  JavaScript, análise de rede e de performance;
  \item \textbf{Git} — o controle de versão que registra o histórico do seu
  código (capítulo~\ref{cap:git});
  \item \textbf{Node.js} — o runtime que executa JavaScript fora do navegador
  (capítulo~\ref{cap:node}).
\end{enumerate}

O apêndice~\ref{apendice:instalacao} traz o passo a passo completo de
instalação e configuração do ambiente, incluindo a verificação de que tudo
funciona (\texttt{node --version}, \texttt{git --version}).

\begin{caixaarmadilha}
Erro comum de iniciante: instalar dezenas de ferramentas antes de precisar
delas. Comece com o mínimo (editor + Git + Node) e instale o resto conforme os
capítulos pedirem. Ambiente simples = foco no aprendizado.
\end{caixaarmadilha}

\section{O full stack na indústria: do júnior ao sênior}

A carreira full stack costuma seguir esta progressão:

\begin{itemize}[leftmargin=*, itemsep=0.4em]
  \item \textbf{Júnior}: executa tarefas bem definidas com supervisão; domina
  o fluxo básico (editar código, rodar, testar, commitar);
  \item \textbf{Pleno}: entrega features completas de ponta a ponta com
  autonomia; escreve testes; participa de code review; entende o sistema como um todo;
  \item \textbf{Sênior}: projeta arquiteturas, lidera decisões técnicas, avalia
  trade-offs (custo, performance, segurança), orienta outros devs e previne
  problemas antes que aconteçam.
\end{itemize}

Este livro é organizado exatamente nessa ordem: os capítulos da Parte~I e~II
formam o nível júnior; as Partes~III e~IV, o nível pleno; a Parte~V (arquitetura,
observabilidade, cloud, IA) e o capítulo de carreira, o nível sênior.

\begin{caixadica}
A Inteligência Artificial mudou o mercado: escrever código está mais fácil, mas
\emph{entender}, \emph{validar} e \emph{arquitetar} está mais valioso. O
diferencial profissional de 2026 é a capacidade de julgar o que a IA produz —
exatamente o que este livro treina com profundidade no capítulo~\ref{cap:ia}.
\end{caixadica}

% ============================================================================
\section{Projeto do capítulo: seu portfólio no GitHub Pages}
\label{sec:projeto1}

\begin{caixaprojeto}
\textbf{Objetivo:} criar sua primeira página profissional — um portfólio
pessoal — e publicá-la na internet de graça, com o GitHub Pages.

\textbf{Critérios de aceite:}
\begin{itemize}[leftmargin=*, itemsep=0.2em]
  \item Página acessível em \texttt{https://seuusuario.github.io};
  \item Contém: título, apresentação, seção de habilidades e contato;
  \item HTML semântico válido (sem erros de validação);
  \item Estilo básico em CSS (cores, espaçamento, fonte);
  \item Publicada via GitHub Pages (branch \texttt{main}, pasta raiz).
\end{itemize}
\end{caixaprojeto}

\subsection{Passo a passo}

\begin{passos}
  \item Crie uma conta gratuita em \url{https://github.com} (se ainda não tiver);
  \item Crie um repositório público chamado \texttt{seuusuario.github.io}
  (exatamente esse formato);
  \item No repositório, adicione os três arquivos abaixo (via \emph{Add file}),
  ou clone o repositório e trabalhe localmente — veja o capítulo~\ref{cap:git};
  \item Salve, aguarde alguns segundos e acesse
  \texttt{https://seuusuario.github.io}.
\end{passos}

\subsection{O código}

O arquivo principal é o \texttt{index.html} (Listagem~\ref{list:portfolio}):

\lstinputlisting[style=livro, language=HTML,
  caption={Arquivo \texttt{codigo/cap01-portfolio/index.html} — estrutura do portfólio.},
  label={list:portfolio}]
  {\codigopath/cap01-portfolio/index.html}

O CSS (\texttt{style.css}, Listagem~\ref{list:portfolio-css}) dá identidade
visual com cores, espaçamento e responsividade básica — conceitos que serão
aprofundados no capítulo~\ref{cap:css}:

\lstinputlisting[style=livro, language=CSS,
  caption={Arquivo \texttt{codigo/cap01-portfolio/style.css} — estilo do portfólio.},
  label={list:portfolio-css}]
  {\codigopath/cap01-portfolio/style.css}

\begin{caixadica}
O projeto completo — incluindo o \texttt{script.js} que anima o contador de
caracteres do formulário e o \texttt{README.md} com instruções — está em
\texttt{codigo/cap01-portfolio/}. Você pode executar localmente com um
servidor simples: \texttt{npx serve codigo/cap01-portfolio}.
\end{caixadica}

\section{Exercícios de fixação}

\begin{exercicios}
  \item Explique, com suas palavras, o que diferencia um desenvolvedor
  frontend, um backend e um full stack.
  \item Liste as três tecnologias mais usadas segundo a pesquisa Stack
  Overflow 2025 e a fonte oficial de cada número.
  \item Descreva o caminho completo de uma requisição HTTP ao acessar um site:
  do DNS até a renderização.
  \item O que o DevTools do navegador mostra na aba \emph{Rede}? Por que isso
  é útil?
  \item Qual foi o marco histórico registrado pelo Octoverse 2025 e por que
  ele importa para a sua escolha de linguagem?
  \item Crie um arquivo \texttt{index.html} mínimo (título, um parágrafo e uma
  imagem) e abra-o no navegador. Valide-o em \url{https://validator.w3.org}.
\end{exercicios}

\section{Desafios}

\begin{desafios}
  \item \textbf{Desafio do repórter:} escolha uma das tecnologias da tabela
  \ref{tab:mercado} e escreva um parágrafo de 5 linhas, com fonte oficial, sobre
  por que ela é líder na sua categoria. Salve como \texttt{pesquisa/pesquisa-pessoal.md}.
  \item \textbf{Desafio do arquiteto:} desenhe (papel ou ferramenta digital) o
  diagrama de uma aplicação full stack típica com pelo menos 5 camadas e
  explique cada seta.
  \item \textbf{Desafio da curiosidade:} abra o DevTools de um site que você
  usa todos os dias, identifique uma requisição que retornou status \texttt{200}
  e outra que retornou \texttt{304}. Explique a diferença.
\end{desafios}

\section{Exercícios de nível sênior}

\begin{seniores}
  \item \textbf{Trade-offs:} compare, em uma tabela, as stacks
  TypeScript/Next.js, Java/Spring e Python/Django para: (a) velocidade de
  desenvolvimento, (b) mercado de vagas no Brasil, (c) ecossistema de IA e
  (d) custo de infraestrutura. Cite pelo menos uma fonte por linha.
  \item \textbf{Leitura crítica:} acesse a pesquisa Stack Overflow 2025 e
  identifique uma métrica que possa \emph{enganar} um leitor apressado (por
  exemplo: ``admirada'' vs.\ ``usada''). Escreva 10 linhas sobre por que
  interpretar dados de mercado exige cuidado.
  \item \textbf{Visão de sistema:} um produto precisa de uma funcionalidade de
  busca. Liste as camadas do sistema (frontend, API, banco, infraestrutura)
  que são afetadas e, para cada uma, um risco técnico possível.
\end{seniores}

\section{Armadilhas comuns}

\begin{itemize}[leftmargin=*, itemsep=0.4em]
  \item \textbf{Aprender dezenas de tecnologias ao mesmo tempo}: escolher
  ``a stack da semana'' antes de dominar fundamentos é o erro nº 1 de iniciantes.
  \item \textbf{Decorar em vez de entender}: memorizar sintaxe sem entender o
  fluxo da requisição HTTP gera profissionais que travam fora do tutorial.
  \item \textbf{Ignorar as fontes}: copiar ``tendências'' de redes sociais sem
  verificar os dados oficiais (Stack Overflow Survey, Octoverse) leva a decisões ruins.
  \item \textbf{Não publicar nada}: conhecimento sem entrega visível (como o
  portfólio deste capítulo) não gera oportunidades de carreira.
\end{itemize}

\section{Resumo do capítulo}

\begin{itemize}[leftmargin=*, itemsep=0.3em]
  \item Full stack é a visão de sistema: frontend, backend, dados e infraestrutura;
  \item Os dados de 2025/2026 (SO Survey, Octoverse, State of JS) apontam para
  JavaScript/TypeScript, React, Next.js, PostgreSQL, Docker e AWS;
  \item Uma requisição HTTP percorre: DNS → conexão segura → servidor → resposta → renderização;
  \item As ferramentas essenciais são editor, terminal, DevTools, Git e Node.js;
  \item O projeto do capítulo — seu portfólio no GitHub Pages — já está
  publicado na internet: primeiro item do seu currículo técnico.
\end{itemize}

\section{Referências oficiais para aprofundar}

\begin{itemize}[leftmargin=*, itemsep=0.3em]
  \item Stack Overflow Developer Survey 2025: \url{https://survey.stackoverflow.co/2025}
  \item GitHub Octoverse 2025: \url{https://octoverse.github.com/}
  \item State of JavaScript 2024: \url{https://2024.stateofjs.com/}
  \item MDN Web Docs (referência definitiva de HTML, CSS e JS): \url{https://developer.mozilla.org}
  \item Roadmap oficial da comunidade full stack: \url{https://roadmap.sh/full-stack}
  \item Validador de HTML do W3C: \url{https://validator.w3.org}
\end{itemize}
